\documentclass[a4paper,10pt]{article}
\usepackage[utf8]{inputenc}
\usepackage[T1]{fontenc}
\usepackage{lmodern}
\usepackage[margin=1.4cm, top=1.2cm, bottom=1.2cm]{geometry}
\usepackage{hyperref}
\usepackage{xcolor}
\usepackage{enumitem}
\usepackage{graphicx}
\usepackage{titlesec}

% Minimal clean styling
\hypersetup{colorlinks=true, urlcolor=blue!70!black, linkcolor=blue!70!black}
\pagestyle{empty}
\setlength{\parindent}{0pt}

% Section formatting
\titleformat{\section}{\large\bfseries\color{black}}{}{0em}{}[\titlerule]
\titlespacing*{\section}{0pt}{6pt}{4pt}

% Colors
\definecolor{accent}{RGB}{20, 50, 120}

% Role command
\newcommand{\role}[4]{%
  \vspace{2pt}\noindent
  \begin{minipage}[t]{0.70\textwidth}\raggedright\textbf{#1} \textbullet\ \textit{#3}\end{minipage}%
  \hfill
  \begin{minipage}[t]{0.28\textwidth}\raggedleft\footnotesize\color{gray}#2 \textbullet\ #4\end{minipage}%
  \vspace{1pt}%
}
\setlist[itemize]{leftmargin=1.2em, itemsep=1.5pt, topsep=2pt, parsep=0pt}

\begin{document}

% ================= HEADER =================
\begin{minipage}[c]{0.80\textwidth}
  {\Huge\bfseries Prabhavit Gupta}\\[3pt]
  {\color{accent}\large AI Engineer $\vert$ Architetture RAG Agentiche Multi-Turno, LlamaIndex, Gemini $\&$ Pgvector}\\[6pt]
  {\footnotesize +39 351 610 7807 ~$\vert$~
  \href{mailto:prabhavitg@gmail.com}{prabhavitg@gmail.com} ~$\vert$~
  \href{https://www.linkedin.com/in/prabhavit/}{linkedin.com/in/prabhavit} ~$\vert$~
  \href{https://github.com/prab-gupta}{github.com/prab-gupta} ~$\vert$~
  Venezia / Milano, Italia (Residente Permanente UE $\vert$ Presenza Flessibile Milano $\vert$ Immediata)}
\end{minipage}\hfill

\vspace{6pt}

% ================= EXECUTIVE SUMMARY =================
\section{Profilo Professionale / Career Summary}
\begin{itemize}
  \item \textbf{Ingegneria di Sistemi AI \& Architetture RAG Agentiche Multi-Turno:} Oltre 3 anni di esperienza nello sviluppo ed industrializzazione di applicazioni AI complesse. Specializzato nella progettazione di architetture RAG conversazionali multi-turno con orchestrazione agentica (LlamaIndex, LangGraph), gestione dello stato di sessione, query transformation dinamica, disambiguazione contestuale e verifica della fedeltà delle risposte.
  \item \textbf{Backend Engineering ad Alte Prestazioni con Python \& FastAPI:} Comprovata esperienza nella costruzione di microservizi asincroni ad alta velocità in Python con FastAPI, pydantic per validazione rigorosa degli schemi JSON e integrazione nativa delle API di modelli di frontiera (Google Gemini 1.5 Pro / Flash, OpenAI, Anthropic Claude).
  \item \textbf{Archiviazione Vettoriale \& Caching Scalabile con PostgreSQL, Pgvector e Redis:} Implementazione di pipeline di retrieval ibrido ad alte prestazioni su PostgreSQL (estensione pgvector per indicizzazione HNSW/IVFFlat), combinate con semantic search, reranking e livelli di caching in-memory su Redis per minimizzare la latenza inferenziale e i costi operativi.
\end{itemize}

% ================= TECHNICAL SKILLS =================
\section{Competenze Tecniche / Technical Skills}
\begin{itemize}[leftmargin=1.2em, itemsep=2pt]
  \item \textbf{Linguaggi \& Framework di Sviluppo:} Python (Avanzato), FastAPI, AsyncIO, PyDantic, TypeScript, Node.js, SQL, Bash.
  \item \textbf{Agentic AI \& LLM Frameworks:} LlamaIndex (Index Query Engines, Sub-Question Querying, Multi-Document Agents), LangGraph, LangChain, Tool Calling, Multi-Turn ReAct Patterns, Google Gemini API, OpenAI API, Anthropic Claude.
  \item \textbf{Database Vettoriali \& Storage:} PostgreSQL, Pgvector (HNSW Indexing, cosine/L2 distance), Redis (Session State, Semantic Cache, Vector Memory), Qdrant, ChromaDB.
  \item \textbf{Strategie RAG Avanzate:} Hybrid Search (Sparse BM25 + Dense Embeddings), Cross-Encoder Reranking (Cohere, BGE), Chunking Semantico, Metadata Filtering, Context Compression, RAGAS Metrics.
  \item \textbf{DevOps \& Messa in Produzione:} Docker, Docker Compose, Git, GitHub Actions CI/CD, Linux Server Management, Logging \& Tracing strutturato, Protezione contro Prompt Injection.
\end{itemize}

% ================= PROFESSIONAL EXPERIENCE =================
\section{Esperienza Professionale / Work Experience}

\role{AI Engineer \& Solutions Architect}{2024 -- Presente}{Indipendente / Consulenza AI Enterprise}{Venezia \& Milano, IT}
\begin{itemize}
  \item Progettazione e implementazione di motori RAG agentici multi-turno end-to-end con LlamaIndex e LangGraph, consentendo agli utenti di interrogare basi di conoscenza complesse attraverso dialoghi continui con conservazione della memoria episodica.
  \item Integrazione di modelli multimodali e long-context di Google Gemini (Gemini 1.5 Pro e Gemini Flash) con prompt strutturati per l'estrazione ad alta precisione da documenti tecnici e contrattuali.
  \item Costruzione del layer di persistenza vettoriale su PostgreSQL con estensione Pgvector, ottimizzando query HNSW per dataset ad elevata dimensionalità e riducendo il tempo di latenza medio a meno di 45ms.
  \item Implementazione di un layer di semantic caching basato su Redis per memorizzare e riutilizzare risposte ed embedding di query frequenti, abbattendo i costi di chiamata API del 38\%.
  \item Sviluppo di API RESTful robuste e documentate con FastAPI e Pydantic, garantendo tipi rigidi, streaming token (SSE) ed autenticazione sicura per l'integrazione con frontend aziendali.
  \item Creazione di test automatici di regressione per prompt e verifica della precisione semantica con benchmark RAGAS per impedire allucinazioni.
\end{itemize}

\vspace{3pt}
\role{Full-Stack AI Software Engineer}{2023 -- 2024}{Tech Automation Systems}{Milano / Remoto, IT}
\begin{itemize}
  \item Sviluppo di microservizi backend in Python (FastAPI) e pipeline asincrone per l'automazione dei processi documentali e l'analisi automatizzata di flussi informativi.
  \item Progettazione di database relazionali su PostgreSQL e gestione di code di messaggi asincrone con Redis per carichi di lavoro batch e data ingestion distribuiti.
  \item Integrazione di modelli linguistici di grandi dimensioni per la classificazione automatica delle richieste, categorizzazione degli intenti ed estrazione strutturata di entità.
  \item Collaborazione attiva in team multidisciplinari (sviluppatori, architetti, PM) adottando metodologie Agile/Scrum, revisione del codice e pratiche di continuous integration.
\end{itemize}

% ================= NOTABLE PROJECTS =================
\section{Progetti di Rilievo / Notable Projects}
\begin{itemize}
  \item \textbf{Motore RAG Conversazionale Multi-Turno (LlamaIndex + Pgvector):} Piattaforma di knowledge retrieval che combina interrogazione gerarchica dei documenti via LlamaIndex, storage vettoriale HNSW su Pgvector e memoria contestuale multi-sessione salvata in Redis.
  \item \textbf{High-Throughput Agentic Gateway (FastAPI + Google Gemini):} Gateway API per l'orchestrazione di agenti con function calling nativo, streaming asincrono in real-time e caching semantico, validato per oltre 100 richieste concorrenti con latenza sub-secondo.
\end{itemize}

% ================= EDUCATION & CERTIFICATIONS =================
\section{Formazione \& Lingue / Education \& Languages}
\begin{itemize}
  \item \textbf{Bachelor of Arts in Philosophy \& Economics} -- Università Ca' Foscari Venezia (2018 -- 2022) \\
  \textit{Formazione avanzata in logica matematica, modelli decisionali, teoria dei giochi ed analisi dei sistemi.}
  \item \textbf{Lingue:} Inglese (C1 Professionale / Fluente), Italiano (Lavorativo / Fluente nei contesti professionali).
  \item \textbf{Stato di Residenza e Fiscale:} Cittadino extra-UE con Permesso di Soggiorno UE a Tempo Indeterminato (pieni diritti di lavoro in Italia/UE senza necessità di sponsorizzazione); Partita IVA attiva / regime ordinario o forfettario pronto.
\end{itemize}

\end{document}
